\section{Scott's semantic framework: explicit arguments and comparison proofs}

\begin{annexitem}{annex:p8.sa.indexed-virtuals}{thm:p08-sa-indexed-virtuals}{Exactly image quantification and the conditional class diagonal}
\begin{proof}
Every value g(x) is in the image, so truth of a predicate on all
of \(g[D]\) implies truth for every indexed argument. Conversely
each \(m\in g[D]\) equals g(x) for some x; truth of every
indexed instance implies truth at m. The existential directions
are the same witness argument. The choice of several x with the
same image m merely repeats one truth value and changes neither
quantifier. In the cardinal illustration the image is precisely
the virtual cardinal classes of the sets quantified over; no
other virtual object is assigned to the cardinal variable.

For the stated class-theoretic diagonal, suppose c=g(a) for
\(a\in D\). Its membership clause gives
\(a\in c\) iff \(a\notin g(a)\), hence iff
\(a\notin c\), contradicting two-valued membership. Therefore
c is outside the image. This inference depends on the allowed
virtual diagonal class and a membership relation on D; without
those assumptions an arbitrary map into an arbitrary V can
certainly be surjective. Nothing here proves an elimination
metatheorem for unspecified virtual-object axioms. It proves
exactly the source's indexed-quantification argument and its
explicit conditional limit.
\end{proof}
\end{annexitem}

\begin{annexitem}{annex:p8.sa.descriptions}{thm:p08-sa-descriptions}{Undefined terms, descriptions and the substitution distinction}
\begin{proof}
At a fixed index i, extend the metalevel value set by a distinguished
undefined tag not in V, solely to express the equality test. A
term's partial value is represented by its V value when defined
and by that tag otherwise. The source equality is exactly ordinary
equality of these tagged values. It is reflexive, symmetric and
transitive in all cases, including two undefined terms; precisely
one undefined term is unequal to a defined one. The tag is not
introduced as a denoted virtual individual.

Let \(B_i=\{b\in D:\|\Phi(\bar b)\|_i=1\}\).
If \(B_i=\{a\}\), the description denotes a. For any
\(y\in D\), its equality to that description is true iff
y=a. The right side \(\forall x(x=y\leftrightarrow\Phi(x))\)
is true iff \(B_i=\{y\}\), also iff y=a. If \(B_i\)
is empty or has more than one member, the description is undefined.
No defined y in D equals it, and no singleton \(\{y\}\)
equals \(B_i\), so both sides of the biconditional are false
for every y. These cases prove the complete description law.
The antecedent of the null-description implication is true
exactly in this latter improper case, and the contradictory
description is always undefined; the source equality then makes
the implication's conclusion true at i. If the original description
is proper, the implication has false antecedent. This explains
why the displayed statement is local, not full concept equality.

For a dotted description, its witness set is
\(A_i\cap B_i\). For a possible-domain individual b, the
formula \(\exists_{\cdot}y(y=\bar b\land\Phi(y))\)
is true at i precisely when b belongs to that same set. Unique
description of this formula over D therefore has exactly the
same defined value, or undefined value, at every i. This proves
the asserted full-concept identity of the two descriptions.

To justify arbitrary replacement of equal rigid constants, if
\(\bar a=\bar b\) holds at i, then a=b, since their values
are defined at all indices. Their entire constant concepts are
equal. Semantic compositionality gives replacement: induct up
the finite expression tree; Boolean operations and equality are
functions of their input concepts, general nonlogical operations
are functions of their input concepts, and at a binder equality
of all argument concepts for each rigid bound-variable value
gives equality of the two whole input functions. Thus the outputs
are equal. For the quantified replacement law, each assignment
of x,y into D supplies two rigid constants. If they are unequal
the antecedent is false; if equal the just established argument
gives equal truth concepts for their formula instances. This
proves the universally quantified implication.

For arbitrary closed terms, use \(I=\{0,1\}\),
\(D=V=\{a,b\}\) with a\(\ne\)b. Let the description
\(\tau=\mathsf I xP(x)\) have singleton predicate
\(P_0=P_1=\{a\}\), and let
\(\sigma=\mathsf I xQ(x)\) have
\(Q_0=\{a\},Q_1=\{b\}\). All these predicates may be
ordinary index-varying extensional ones. At 0 the terms are equal.
Take \(\Phi(z)=\Box(z=\bar a)\). Then \(\Phi(\tau)\)
is true and \(\Phi(\sigma)\) false at 0. Local equality
cannot justify replacement. In contrast,
\(\tau\equiv\sigma\) states equality of the complete
partial functions, with equal domains of definition as well as
equal values. The same compositional induction gives arbitrary
capture-free replacement of such closed terms.
\end{proof}
\end{annexitem}

\begin{annexitem}{annex:p8.sa.extensionality}{thm:p08-sa-extensionality}{Extensional predicates and replacement of whole formula concepts}
\begin{proof}
Suppose \(\tau=\tau'\) and \(\sigma=\sigma'\) at i.
If \(\tau\bar R\sigma\) is true, both its arguments
are defined, so the equal replacement arguments are defined
with the same values. Membership in the ordinary relation is
unchanged. If that predicate is false, either an argument is
undefined or the pair is outside the relation. The matched
definedness and values leave it false after replacement as well.
The proof is identical for the varying \(R_i\). The source
equality is locally extensional by its tagged-value interpretation,
but at two undefined arguments it is true. No strict ordinary
binary predicate can have this value there. Since the language
has the contradictory description, those argument concepts do
occur, so this is an actual separation, not an empty case.
For the source pullback comparison, membership of (a,b) in
\(R_i\) is by definition membership of
\((t_i(a),t_i(b))\) in \(R_0\); therefore the two atomic
evaluations agree. Nothing about this equality of truth tests
requires \(t_i\) to be injective. If it is not, equality of
the transformed individuals can hold for distinct a,b, so a
full preservation of rigid individual equality would be an
unjustified additional claim.

For a general predicate, equality of its whole input concepts
means that the very same arguments enter its fixed total function.
Its output concept is therefore equal. The same argument applies
to every general propositional operation and to any prescribed
finite arity. For closed formulas,
\(\Box(\Phi\leftrightarrow\Psi)\) at any index states
that their truth concepts are equal on all I. Induct up each
chosen occurrence: Boolean parents preserve equality, general
operation parents preserve equality by their function types,
and binders preserve it when the whole rigid-instance argument
functions agree. Capture-free closed replacement does not alter
the value supplied to any bound variable. Thus all resulting
truth concepts agree. For example, p and q can agree at index 0
but not at 1; then \(\Box p\) and \(\Box q\) can disagree
already at 0. Agreement of current extensions is weaker than
the whole-concept hypothesis actually used.
For the logical-operation warning, let a permutation \(\pi\) act
on truth concepts by \((\pi f)(i)=f(\pi^{-1}(i))\).
The constant-1 map F satisfies \(F(\pi f)=1_I=\pi F(f)\)
for every f and \(\pi\), and the same equation with \(0_I\)
holds for the constant-0 map. Both commute with all permutations.
They are not negation, double negation, universal box or its dual
on a nonempty index domain. Thus that invariance criterion alone
cannot justify the source's suggested four-operation conclusion.
\end{proof}
\end{annexitem}

\begin{annexitem}{annex:p8.sa.universal-quantifiers}{thm:p08-sa-universal-quantifiers}{Universal S5 soundness and both actual-domain countermodels}
\begin{proof}
The universal box of a truth concept f is the constant-1 concept
when every f(j)=1, and the constant-0 concept otherwise. Dual
diamond therefore tests whether f(j)=1 for some j. If both
\(\Box(\Phi\to\Psi)\) and \(\Box\Phi\) are true,
every index has both inner antecedents true, so every index
satisfies \(\Psi\); this proves K. Box truth implies truth
at the current index, proving T. A true box is true at all
indices, proving 4. A false box is false at all indices, so
its negation is everywhere true and universally necessary;
this proves 5. Duality follows by Boolean complement. If a
formula is logically valid across all interpretations and
indices, every index of each particular interpretation satisfies
it, so its box is valid too. Mere truth at one chosen index
is plainly not that premise.

For the possible-domain commuting law, the left side says
\(\|\Phi(\bar a)\|_j=1\) for every j in I and every
a in D. The right side says the same condition with the two
universal quantifiers in the opposite order. D is fixed over
all indices, so these are equivalent even for general intensional
\(\Phi\). Actual restriction at i follows because a rigid
\(\bar a\) satisfies \(\widetilde A\) precisely when
\(a\in A_i\); implication ignores all other a.

For failure from left to right of the actual commuting law,
let \(I=\{0,1\}\), \(D=V=\{a,b\}\),
\(A_0=\{a\},A_1=\{b\}\), and let a strict unary
predicate have \(P_0=\{a\},P_1=\{b\}\).
At every index every actual individual satisfies P, so
\(\Box\forall_{\cdot}xP(x)\) is true at 0.
But a is actual at 0 and fails P at 1; hence
\(\forall_{\cdot}x\Box P(x)\) is false at 0.
For failure from right to left take instead
\(A_0=\{a\},A_1=\{a,b\}\),
\(P_0=P_1=\{a\}\). At 0 every currently actual a
satisfies P at every index, so the right side is true. At 1
the actual b fails P, so the universal left side is false.
All actual domains in these examples are nonempty; allowing
empty ones is not needed to produce the separation.
\end{proof}
\end{annexitem}

\begin{annexitem}{annex:p8.sa.neighborhood-bridge}{thm:p08-sa-neighborhood-bridge}{Full representation and the exact later-system comparison}
\begin{proof}
Each truth concept \(f:I\to2\) equals the characteristic
function \(1_X\) of \(X=\{i:f(i)=1\}\). Thus
\(F(f)(i)=1\) iff that X is in the stated
\(\mathscr P_i\). Conversely given arbitrary families,
this equation defines a unique total F, since it assigns 0
or 1 to every input concept at every i. Recovering the families
from this F returns exactly their original memberships. These
two constructions are inverse, so the representation is unique.

Use the same index/world set and the same atomic truth sets.
Induct on the common propositional language. Atoms agree by
definition, Boolean connectives by their pointwise clauses,
and the modal case by the exact membership equivalence just
proved and the inner truth-set induction. This gives truth
preservation at every index, including the defined dual operator.
Global validity of an inner biconditional in an interpretation
means equality of its truth sets, hence equality of modal outputs.
The finite two-index example with \(p_0=q_0=1,p_1=1,q_1=0\)
and \(N(0)=N(1)=\{I\}\) falsifies lifting only current
equivalence: at 0 p and q agree, but their boxes do not.

For the editorial M/C/N comparison, let X,Y be formula truth sets.
Upward closure accepts each factor when it accepts X\(\cap\)Y,
giving M; intersection closure accepts X\(\cap\)Y when it
accepts both, giving C; inclusion of I gives boxed truth, N.
Conversely a supplementation failure supplies
\(X\subseteq Y,X\in N(i),Y\notin N(i)\); assign two atoms
truth sets X,Y, so their conjunction has truth set X and M fails.
An intersection failure is witnessed by two atoms with its
accepted factors; N fails exactly when its constant truth set
I is absent. These are all-valuations converses, not conditions
inferred from one fixed valuation.

To distinguish unit from normality, on \(I=\{0,1,2\}\)
put \(N(i)=\{\{0\},I\}\) for all i, and assign
\(T(p)=\{0\},T(q)=\{0,1\}\). Both \(\Box p\)
and \(\Box(p\to q)\) are true, while \(\Box q\) is
false. This unit-containing frame fails K. On one world set
\(N(i)=\varnothing\); supplementation and intersections
hold vacuously, but N fails. This separates EMC from EMCN.
The later naming and completeness identifications depend on
the exact common defined-dual language, propositional base,
theorem-only modal rules and local finite-premise deduction
proved in the Chellas and Pacuit presentations. They are not
asserted to be Scott's 1970 theorem inventory. Scott's universal
box accepts precisely the one truth set I, giving
\(\mathscr P_i=\{I\}\) as stated.
\end{proof}
\end{annexitem}

\begin{annexitem}{annex:p8.sa.progressive}{thm:p08-sa-progressive}{Interior soundness and failure of same-domain relational representation}
\begin{proof}
For a truth set X, the progressive clause says that i lies in
an open interval contained in X, precisely \(i\in\operatorname{Int}(X)\).
Interior contains the whole space when X does, is monotonic,
preserves binary intersections, and satisfies
\(\operatorname{Int}(X)\subseteq X\).
The first three properties give the normal K principles;
the inclusion gives T. Also \(\operatorname{Int}(X)\) is
open, so taking its interior again changes nothing, giving 4.
The dual modality is complement of interior of the complement.
These establish S4 soundness for the specified operator; they
do not prove its full exact-S4 completeness classification.

At a fixed real i, its accepted family consists of sets with i
in their interior. Every member contains i. For any \(j\ne i\),
choose a sufficiently small open interval centered at i omitting
j. The intersection of all accepted sets is therefore exactly
\(\{i\}\). That singleton has empty real-line interior,
so is not accepted. If a relation on this same domain represented
the operator on all subsets, its successor set \(R(i)\)
would be the intersection of all its accepted supersets and
would itself be accepted. This contradicts the computed family.
Arbitrary atomic valuations supply all subsets as arguments, so
the obstruction is to the required operation equality, not merely
to one restricted formula valuation. A relational semantics on
a different domain can have the same S4 theorem set without
being this pointwise operation on \(\mathbb R\).
\end{proof}
\end{annexitem}

\begin{annexitem}{annex:p8.sa.binder-replacement}{thm:p08-sa-binder-replacement}{Whole argument functions and binding replacement}
\begin{proof}
At a chosen index suppose the two hypotheses hold. For each
\(a\in D\), the first gives
\(\|\tau(\bar a)\|=\|\sigma(\bar a)\|\) as partial
functions on all I, with the same definedness at every point.
Thus the two functions \(D\to V^{(I)}\) supplied to the
binder are equal. The second hypothesis similarly makes
\(\|\Phi(\bar a)\|=\|\Psi(\bar a)\|\) on all I
for each a, so the second two binder arguments agree as
\(D\to2^I\) functions. The prescribed binder interpretation
is a function of this pair, hence has the same output concept.
For a term maker, equality of those complete partial outputs
is exactly the asserted intensional equality. For a formula
maker it is exactly the universal boxed biconditional.

To see why equality only of current extensions cannot replace
the hypothesis, take a singleton D, two indices 0,1, and two
formula argument functions which agree at 0 but differ at 1.
A type-correct formula binder may return the constant truth
concept testing whether the sole input formula concept is true
at both indices. Its outputs differ already at 0. A term binder
can analogously return one of two distinct rigid individuals
according to that test. These permitted total binder functions
show that pointwise replacement is not sound in the general type.
They do not assume the source's unproved proposal that every
binder has a reduction to its simpler logical binders.
\end{proof}
\end{annexitem}

\begin{annexitem}{annex:p8.sa.incidence}{thm:p08-sa-incidence}{The two function layers and finite witnesses to every failed converse}
\begin{proof}
Under \(V\subseteq S^I\), rigid a,b are equal iff their
whole history functions agree, iff \(a_j=b_j\) at every j.
That last condition is exactly \(\Box(\bar a\approx\bar b)\).
For arbitrary closed terms, full concept equality implies
equality of denoted individuals at every index, including matched
undefinedness. At each index, equal defined individuals have
the same current state; two undefined concepts satisfy both
local equality and incidence by convention. Thus universal
equality implies universal incidence, and each universal
condition implies its own current-index condition. This proves
all displayed arrows. Whole concept equality is universal local
equality by its definition, so the initial arrow is reversible.

For concrete failures use \(I=\{0,1\}\), \(S=\{0,1\}\)
and distinct total histories
\(a=(0,0),b=(0,1),c=(1,0),d=(1,1)\) in V=D.
Let \(\tau\) have constant concept a. Let \(\sigma\)
denote a at 0 and b at 1. At 0 local equality holds. At 1
the current states are \(a_1=0,b_1=1\), so incidence fails;
\(\Box(\tau\approx\sigma)\) is false. Thus local equality
does not imply universal incidence or full concept equality.

Now let \(\sigma\) denote b at 0 and c at 1. Its diagonal
states are \(b_0=0,c_1=0\), equal to those of the constant
a concept at both indices. Universal incidence therefore holds,
but the denoted histories differ at both indices, so local
equality fails everywhere and \(\Box(\tau\ne\sigma)\)
is true. This proves that universal incidence does not imply
local equality, universal equality or full concept equality.
Already the two rigid histories a and b are incident at 0 but
distinct, showing failure of the local incidence-to-equality
arrow. The first constructed example shows that current
incidence also does not imply its universal version. The varying
concepts can be realized as unique descriptions of strict unary
predicates selecting their respective history at each index,
so they are available within the source language.

If instead states reflect a restricted nonseparating property
family, take distinct individuals a,b and assign both the same
state at every index. The constant names then have universal
incidence but are unequal. This is not a counterexample under
\(V\subseteq S^I\), which would identify those identical
histories. It shows exactly why dropping the separation
hypothesis, as contemplated in the source discussion, invalidates
the reconstruction of individual equality. The source highway
and organism examples concern these different choices of what
the states individuate; the hypotheses cannot be silently merged.
\end{proof}
\end{annexitem}

\begin{annexitem}{annex:p8.sa.counterparts}{thm:p08-sa-counterparts}{Counterpart profiles and the source's one-direction implication}
\begin{proof}
At a fixed i define \(H_i(a)=\{z\in A_i:C_i(z,a)\}\).
The counterpart definition of \(a\approx b\) is exactly
\(H_i(a)=H_i(b)\). Equality of these sets is reflexive,
symmetric and transitive, so profile incidence is an equivalence
relation. If a is actual, the required reflexive pair
\(C_i(a,a)\) puts a in \(H_i(a)\). If a\(\approx\)b,
it is also in \(H_i(b)\), giving \(C_i(a,b)\).
This proves the source implication for every D-assignment.

For its failed converse let \(I=\{0,1\}\),
\(A_0=\{a,z\},A_1=\{b\}\), and
\(D=A_0\cup A_1\). Define
\(C_0=\{(a,a),(z,z),(a,b),(z,b)\}\) and
\(C_1=\{(b,b)\}\). Every first component belongs to the
current actual set; a pair with current actual second component
is an identity pair; and every current actual has its reflexive
pair. Thus all three required conditions hold in both worlds.
At 0, \(C_0(a,b)\) is true, but
\(H_0(a)=\{a\}\ne\{a,z\}=H_0(b)\). Incidence is
false there. This is a full two-index model of the source
conditions, not an unspecified counterpart example.
\end{proof}
\end{annexitem}
